\documentclass[11pt]{article}

\usepackage[margin=1in]{geometry}
\usepackage{amsmath, amssymb, amsthm}
\usepackage{mathtools}
\usepackage{bm}
\usepackage{hyperref}
\usepackage{enumitem}
\usepackage{tikz}
\usepackage{microtype}

\title{Multiplicity-Theoretic Framing of the Quantum Prime Abacus, PME, and PL-Commit}
\author{Draft for Internal Use}
\date{\today}

\hypersetup{
  colorlinks=true,
  linkcolor=blue,
  citecolor=blue,
  urlcolor=blue
}

\begin{document}
\maketitle

\begin{abstract}
This report reinterprets the Quantum Prime Abacus (QPA), Prime-Modular Hypergraph Partitioning (PMHP), Prime-Multiplicity Extraction (PME), and the PL-Commit commitment scheme within the language of Multiplicity Theory and the Genius v2 ``prime moves'' framework. We provide: (i) an executive-level summary for cryptographic and quantum-engineering stakeholders, (ii) a comprehensive mathematical overview of the QPA architecture and associated hardness assumptions, (iii) a canonical alphabet of prime moves and an explicit multiplicity-calculus ``program'' for encoding PMHP into QPA, (iv) a dynamics + shadows + inversion phase written in the same calculus, and (v) code-style snippets that illustrate how these ideas could be exposed in a software stack such as PhaseMirror-HQ.
\end{abstract}

\tableofcontents

\section{Executive Summary}

\subsection{Conceptual Goal}

The overarching goal is to treat the Quantum Prime Abacus (QPA) not merely as a quantum algorithmic gadget, but as a concrete instance of a prime-indexed multiplicity architecture. In Multiplicity Theory, mathematical and computational structures are understood as recursively generated patterns of prime-labeled interactions, with identity and behavior preserved through feedback across scales. The QPA instantiates this view as a three-layer quantum system:

\begin{enumerate}[label=(\alph*)]
  \item A \emph{physical layer} of qubits or qudits (hardware-agnostic).
  \item A \emph{prime-bead logical layer} where each bead carries a prime label, a module/partition index, and a local multiplicity register.
  \item A \emph{multiplicity semantics layer} where partition structure is encoded in recursive prime factorization patterns.
\end{enumerate}

On top of this architecture, the PMHP problem is introduced as a prime-labeled hypergraph partitioning task with modular arithmetic constraints, and PME is defined as the problem of recovering an optimal or near-optimal partition from classical shadows of an optimal QPA state. These pieces are then combined to construct PL-Commit, a quantum-native cryptographic commitment scheme whose security is argued to rest on the hardness of extracting prime-multiplicity trajectories from entanglement.

\subsection{Why This Matters}

From a cryptographic and systems perspective, the key claims are:

\begin{itemize}
  \item \textbf{New hardness source.} PME is conjectured to be hard for both classical and quantum adversaries due to a combination of computational (NP-hardness), information-theoretic (shadow tomography limits, no-cloning), and learning-theoretic (quantum PAC lower bounds) barriers.
  \item \textbf{Quantum-native security.} PL-Commit does not rely on traditional algebraic trapdoors (e.g., lattices, codes). Instead, security is tied directly to the structure of entanglement and the difficulty of recovering multiplicity trajectories from classical shadows.
  \item \textbf{NISQ-oriented engineering.} The architecture includes an explicit engineering blueprint: Hierarchical Amplitude-Compressed QPA (HAC-QPA), a hybrid quantum-classical decision boundary functional, and a QPACT compiler design.
\end{itemize}

From a Multiplicity Theory perspective, these constructions validate the idea that primes can act as ``eigenmodes'' of computation and that prime-indexed multiplicity spaces provide a natural basis for modeling non-linear, emergent structure in quantum algorithms and cryptography.

\section{Multiplicity-Theoretic Overview of QPA}

\subsection{Prime-Bead Logical Layer}

We reinterpret the prime-bead as a local multiplicity space.

\begin{definition}[Prime-Bead as Multiplicity Space]
Fix a set of primes $\mathcal{P}$ and an interaction hypergraph $G=(V,E)$ with $|V| = N$. For each vertex $v_k \in V$ we associate:
\begin{enumerate}[label=(\roman*)]
  \item A prime label $p_k \in \mathcal{P}$ (a static type-level identity).
  \item A partition index $i \in [K] := \{1,\dots,K\}$ (module label).
  \item A multiplicity register value $\mu_k$ in
  \[
    M_k := \{0,1,\dots,\deg(v_k)\},
  \]
  counting incident hyperedges that are currently ``uncut'' under a partition.
\end{enumerate}
The logical Hilbert space for bead $k$ is then
\[
  \mathcal{H}_k \cong \mathbb{C}^K \otimes \mathbb{C}^{|M_k|},
\]
with basis vectors $\{|i,\mu\rangle_k : i \in [K], \mu \in M_k\}$.
\end{definition}

\begin{definition}[Global Prime-Indexed Multiplicity Space]
The global QPA logical space is
\[
  \mathcal{H} = \bigotimes_{k=1}^N \mathcal{H}_k.
\]
We interpret $\mathcal{H}$ as a prime-indexed multiplicity space: each basis configuration $\bigotimes_k |i_k,\mu_k\rangle_k$ describes how each prime $p_k$ participates in global module structure and local recurrence.
\end{definition}

\begin{definition}[Prime-Multiplicity Potential]
For each bead $k$, define the prime-multiplicity observable
\[
  \hat{P}_k |i,\mu\rangle_k
    := (\log p_k + \alpha \mu)\,|i,\mu\rangle_k,
\]
where $\alpha > 0$ is a coupling constant. The term $\log p_k$ encodes the intrinsic scale of the prime eigenmode at site $k$, and $\mu$ encodes local recurrence in the interaction graph.
\end{definition}

In Multiplicity Theory terms, $\hat{P}_k$ can be read as the local generator of ``identity through recurrence'': it combines prime identity and local multiplicity into a single observable.

\subsection{Recursive Prime-Multiplicity Semantics}

We now capture the semantics of a partition $\sigma:V \to [K]$ in a prime-recursive invariant.

\begin{definition}[Prime Multiplicity Function]
Fix an ordering $V = \{v_1,\dots,v_N\}$ and a partition $\sigma:V \to [K]$. The prime multiplicity function $f_P(\sigma,k)$ is defined recursively by
\[
  f_P(\sigma,1) := p_1,
\]
and for $k \ge 2$,
\[
  f_P(\sigma,k) :=
  \begin{cases}
    \displaystyle
    \dfrac{\gcd\big(f_P(\sigma,k-1),p_k\big)\cdot
           \operatorname{lcm}\big(f_P(\sigma,k-1),p_k\big)}
          {f_P(\sigma,k-1)}
      & \text{if } \sigma(k) = \sigma(k-1),\\[8pt]
    f_P(\sigma,k-1)\cdot p_k
      & \text{if } \sigma(k) \neq \sigma(k-1).
  \end{cases}
\]
\end{definition}

\paragraph{Multiplicity-Theoretic Interpretation.}

At each step $k$, the recursion applies a small prime-indexed move:

\begin{itemize}
  \item If $\sigma(k) \neq \sigma(k-1)$, the system executes a \emph{prime insertion} move: a new prime factor $p_k$ is appended to the trajectory via multiplication.
  \item If $\sigma(k) = \sigma(k-1)$, the system executes a \emph{gcd/lcm closure} move: the existing prime pattern $f_P(\sigma,k-1)$ is refined by intersecting and unifying with $p_k$ via $\gcd$ and $\operatorname{lcm}$, normalized by the previous state.
\end{itemize}

Thus $f_P(\sigma,k)$ encodes the partition path as a prime-factorization invariant, built from an ordered sequence of microscopic prime moves.

\subsection{Prime Couplings and Coherence}

We next capture how prime-labeled sites interact across modules.

\begin{definition}[Prime Coupling Tensor]
For beads $i,j$ with primes $p_i,p_j$ and modules $a,b\in[K]$, define
\[
  T_{ij}^{ab} :=
  \begin{cases}
    -\log\big(\gcd(p_i,p_j)\big), & a = b,\\[4pt]
    +\log\big(\operatorname{lcm}(p_i,p_j)\big), & a \neq b.
  \end{cases}
\]
\end{definition}

\begin{definition}[Two-Bead Coupling Hamiltonian]
The two-bead Hamiltonian $H_{ij}$ is
\[
  H_{ij} := \sum_{a,b\in[K]} T_{ij}^{ab}
    \bigl( |a\rangle\langle a|_i \otimes |b\rangle\langle b|_j \bigr).
\]
\end{definition}

\paragraph{Interpretation.}

\begin{itemize}
  \item When $a=b$ (same module), large $\gcd(p_i,p_j)$ lowers the energy; primes that share factorization structure ``prefer'' to co-occur in the same module.
  \item When $a\neq b$ (different modules), large $\operatorname{lcm}(p_i,p_j)$ raises the energy; widely spread multiplicity patterns are penalized when split across modules.
\end{itemize}

This defines a prime-aware coherence rule: the energy landscape encourages configurations whose module assignments respect prime recurrence structure.

\section{Prime-Modular Hypergraph Partitioning (PMHP)}

\subsection{Problem Definition}

\begin{definition}[Prime-Modular Hypergraph Partitioning (PMHP)]
A PMHP instance consists of:
\begin{itemize}
  \item A hypergraph $G = (V,E)$ with $|V| = N$ nodes and $|E| = M$ hyperedges.
  \item Prime labels $p = (p_1,\dots,p_N) \in \mathcal{P}^N$.
  \item A partition count $K \in \mathbb{N}$.
  \item Modular constraints $\{(d_i,r_i)\}_{i=1}^K$ with $d_i \in \mathbb{N}$ and $r_i \in \mathbb{Z}_{d_i}$.
  \item Edge weights $w = (w_1,\dots,w_M)$.
\end{itemize}
The task is to find a partition $\sigma:V \to [K]$ minimizing
\[
  C(\sigma) := \sum_{e_j\in E} w_j\,\mathbf{1}[\operatorname{cut}(e_j,\sigma)]
\]
subject to
\[
  \prod_{v_k \in \sigma^{-1}(i)} p_k \equiv r_i \pmod{d_i}
  \quad\text{for all } i \in [K].
\]
\end{definition}

The cost measures how many hyperedges are cut by the partition, while modular constraints enforce that each module $i$ has a prescribed aggregate prime-multiplicity footprint (a residue class modulo $d_i$).

\subsection{Hamiltonian Encoding}

The PMHP instance is encoded into a QPA Hamiltonian of the form
\[
  H_C = H_{\text{cut}} + \lambda_{\text{mod}} H_{\text{mod}} + \lambda_{\text{bal}} H_{\text{bal}},
\]
where:

\begin{itemize}
  \item $H_{\text{cut}}$ penalizes cut hyperedges using projectors onto equal module assignments.
  \item $H_{\text{mod}}$ penalizes violations of modular constraints via precomputed weights $\omega(p_k,d_i,r_i)$.
  \item $H_{\text{bal}}$ enforces balanced module sizes by penalizing deviations from $N/K$ nodes per module.
\end{itemize}

This QPA cost Hamiltonian defines an energy landscape on $\mathcal{H}$ whose ground states correspond to near-optimal PMHP solutions.

\section{Prime-Multiplicity Extraction (PME)}

\subsection{Problem Definition}

\begin{definition}[Prime-Multiplicity Extraction (PME)]
Given:
\begin{itemize}
  \item A PMHP instance $(G,p,K,\{(d_i,r_i)\},w)$.
  \item Classical shadows $S = \{(U_j,z_j)\}_{j=1}^m$ of an optimal QPA state $|\Psi^*\rangle$, prepared under $H_C$.
  \item Full specification of $H_C$ and the coupling tensors $\{T_{ij}\}$.
\end{itemize}
The PME problem is to output a partition $\sigma^*:V\to[K]$ such that
\[
  \langle \sigma^*|H_C|\sigma^*\rangle
    \le \langle \Psi^*|H_C|\Psi^*\rangle + \varepsilon,
\]
for small $\varepsilon > 0$, with all modular constraints satisfied.
\end{definition}

\subsection{Multiplicity-Theoretic View}

From a multiplicity perspective, QPA defines a forward map
\[
  \sigma \;\mapsto\; \text{prime-multiplicity trajectory } f_P(\sigma,\cdot)
  \;\mapsto\; \text{entangled state } |\Psi\rangle
  \;\mapsto\; \text{shadows } S,
\]
where:

\begin{itemize}
  \item The first arrow is the recursive construction of $f_P$ from $\sigma$ via prime insertion and gcd/lcm closure.
  \item The second arrow is concentration of amplitude on low-energy trajectories by QPA dynamics under $H_C$.
  \item The third arrow is the classical-shadow measurement process.
\end{itemize}

PME asks to invert this pipeline: given $S$ and the forward rules, recover a partition $\sigma^*$ whose multiplicity trajectory is consistent with the shadows and near-optimal under $H_C$.

\section{PL-Commit as a Multiplicity Commitment}

\subsection{Protocol Summary}

PL-Commit is a commitment scheme built on PME hardness. The setup fixes a \emph{reference} multiplicity landscape, and commitment/opening correspond to choosing and revealing a branch of the multiplicity recursion.

\subsubsection*{Setup}

\begin{enumerate}[label=(S\arabic*)]
  \item Sample a random PMHP instance $I = (G,p,K,\{(d_i,r_i)\})$.
  \item Run QPA under $H_C$ to obtain an optimal (or near-optimal) state $|\Psi^*\rangle$.
  \item Generate $m$ classical shadows $S = \{(U_j,z_j)\}_{j=1}^m$ of $|\Psi^*\rangle$.
  \item Publish $(I,S)$ as public parameters.
\end{enumerate}

\subsubsection*{Commit}

To commit to a message $m \in \{0,1\}^\ell$:

\begin{enumerate}[label=(C\arabic*)]
  \item Encode $m$ as a perturbation of the module-level multiplicity invariants:
  \[
    r_i' := r_i \oplus (m \bmod d_i).
  \]
  \item Run QPA under the perturbed constraints $\{(d_i,r_i')\}$ to obtain the perturbed state $|\Psi_m\rangle$.
  \item Generate shadows $S_m = \{(U'_j,z'_j)\}_{j=1}^{m'}$ of $|\Psi_m\rangle$.
  \item Choose random $r$ and compute a hash $h := \mathrm{hash}(m,r)$.
  \item Output commitment $C = (S_m,h)$.
\end{enumerate}

\subsubsection*{Open}

To open a commitment $C$ to $m$:

\begin{enumerate}[label=(O\arabic*)]
  \item Reveal $(m,r)$, and let the verifier check $h(m,r) = C.h$.
  \item The verifier uses $S_m$ and the known forward rules to reconstruct a partition $\sigma_m$.
  \item The verifier accepts iff $\sigma_m$ satisfies the perturbed modular constraints $\{(d_i,r_i')\}$.
\end{enumerate}

\subsection{Security Intuition}

\paragraph{Binding.} To break binding, an adversary must open the same $C$ to two different messages $m \neq m'$. This would require recovering two distinct partitions $\sigma_m$ and $\sigma_{m'}$ consistent with the same shadow data and simultaneously satisfying two distinct sets of modular constraints. Conceptually, this forces the adversary to solve PME twice on the same shadows, contradicting the assumed hardness of PME.

\paragraph{Hiding.} Hiding is supported by information-theoretic arguments: classical shadows carry only $O(m \log n)$ bits of information about a $2^n$-dimensional quantum state. The message $m$ is encoded as a small perturbation of modular constraints, and $m$ is hidden inside a high-dimensional entangled prime-multiplicity pattern; no efficient reconstruction of $m$ from the shadows alone is known.

\section{Canonical Prime-Move Alphabet}

We now compress the multiplicity-theoretic operations into a small canonical alphabet of ``prime moves'', designed both as a cognitive vocabulary (Genius v2) and as a basis for a multiplicity calculus.

\subsection{Seven Prime Moves}

We introduce seven moves:

\begin{description}
  \item[P-Type (Prime Typing)] Instantiate a local multiplicity space at a site: assign a prime label $p_k$, a module/partition coordinate space $\mathbb{C}^K$, and a multiplicity spectrum $M_k$.

  \item[P-Count (Multiplicity Counting)] Track local recurrence by maintaining $\mu_k \in M_k$, the count of incident hyperedges that remain uncut under a given partition.

  \item[P-Extend (Prime Insertion / Extension)] Extend a running multiplicity trajectory by multiplying in a new prime factor when the module identity changes, e.g.\ $f_P(\sigma,k) := f_P(\sigma,k-1)\cdot p_k$ when $\sigma(k)\neq\sigma(k-1)$.

  \item[P-Close (gcd/lcm Closure)] Refine a running multiplicity invariant by applying gcd/lcm when the module identity persists:
  \[
    f_P(\sigma,k) := \frac{\gcd(f_P(\sigma,k-1),p_k)\cdot
                          \operatorname{lcm}(f_P(\sigma,k-1),p_k)}
                         {f_P(\sigma,k-1)}.
  \]

  \item[P-Cohere (Coherence Scoring / Coupling)] Evaluate and enforce compatibility between two prime sites across module assignments using $T_{ij}^{ab}$: reward same-module placement for primes with large gcd, penalize cross-module placement with large lcm.

  \item[P-Shadow (Shadow Projection)] Project an entangled prime-multiplicity state into classical data by sampling random Clifford measurements and recording outcomes $(U_j,z_j)$.

  \item[P-Invert (Multiplicity Inversion Attempt)] Given classical shadows and knowledge of the forward rules, attempt to reconstruct a hidden prime-multiplicity trajectory (a partition $\sigma^*$) consistent with those shadows and the constraints of $H_C$.
\end{description}

\section{A Tiny Multiplicity-Calculus Program}

We now write the PMHP $\to$ QPA encoding and the subsequent dynamics + shadow + inversion phases as explicit finite sequences of these prime moves.

\subsection{Program: PMHP $\to$ QPA Encoding}

Let $(G,p,K,\{(d_i,r_i)\},w)$ be a PMHP instance.

\subsubsection*{Step 0: Initialize Global Space}

For each vertex $v_k$:

\begin{enumerate}[label=(0.\arabic*)]
  \item \texttt{P-Type(k, p\_k, K)}: instantiate bead $b_k$.
  \item \texttt{P-Count(k, deg(v\_k))}: define multiplicity spectrum $M_k$ and register $\mu_k$.
\end{enumerate}

This yields
\[
  \mathcal{H} = \bigotimes_{k=1}^N \mathbb{C}^K \otimes \mathbb{C}^{|M_k|}.
\]

\subsubsection*{Step 1: Local Prime Potentials}

For each bead $k$:

\begin{enumerate}[label=(1.\arabic*)]
  \item Use \texttt{P-Type} and \texttt{P-Count} to define $\hat{P}_k|i,\mu\rangle_k
    = (\log p_k + \alpha \mu)|i,\mu\rangle_k$.
\end{enumerate}

\subsubsection*{Step 2: Recursive Semantics (Definition of $f_P$)}

Fix an ordering $(v_1,\dots,v_N)$.

\begin{enumerate}[label=(2.\arabic*)]
  \item \texttt{P-Extend(1)}: set $f_P(\sigma,1) := p_1$.
  \item For $k=2,\dots,N$:
  \begin{itemize}
    \item If $\sigma(k)\neq\sigma(k-1)$, execute \texttt{P-Extend(k)}:
    \[
      f_P(\sigma,k) := f_P(\sigma,k-1)\cdot p_k.
    \]
    \item If $\sigma(k)=\sigma(k-1)$, execute \texttt{P-Close(k)}:
    \[
      f_P(\sigma,k) := \frac{\gcd(f_P(\sigma,k-1),p_k)\cdot
                             \operatorname{lcm}(f_P(\sigma,k-1),p_k)}
                            {f_P(\sigma,k-1)}.
    \]
  \end{itemize}
\end{enumerate}

This defines the semantic mapping $\sigma \mapsto f_P(\sigma,k)$.

\subsubsection*{Step 3: Edge Interactions}

For each pair of sites $i,j$ and modules $a,b$:

\begin{enumerate}[label=(3.\arabic*)]
  \item \texttt{P-Cohere(i,j,a,b)}: compute
  \[
    T_{ij}^{ab} :=
    \begin{cases}
      -\log(\gcd(p_i,p_j)), & a = b,\\[3pt]
      +\log(\operatorname{lcm}(p_i,p_j)), & a \neq b.
    \end{cases}
  \]
  \item Assemble $H_{ij} = \sum_{a,b} T_{ij}^{ab} |a\rangle\langle a|_i \otimes |b\rangle\langle b|_j$.
\end{enumerate}

Then build $H_{\text{cut}}$ by combining $H_{ij}$ across hyperedges so that configurations with cut edges incur higher energy.

\subsubsection*{Step 4: Modular and Balance Constraints}

For modular constraints:

\begin{enumerate}[label=(4.\arabic*)]
  \item For each module $i$ and site $k$:
  \begin{itemize}
    \item Use \texttt{P-Extend}/\texttt{P-Close} at the module-invariant level to define weights $\omega(p_k,d_i,r_i)$.
    \item Contribute $\omega(p_k,d_i,r_i) |i\rangle\langle i|_k$ to $H_{\text{mod}}$.
  \end{itemize}
\end{enumerate}

For balance:

\begin{enumerate}[label=(4.\arabic*')]
  \item Use \texttt{P-Count} at the module level to encode
  \[
    H_{\text{bal}} := \sum_{i=1}^K \left(\sum_{k=1}^N |i\rangle\langle i|_k - \frac{N}{K}\right)^2.
  \]
\end{enumerate}

\subsubsection*{Step 5: Final QPA Cost Hamiltonian}

\begin{enumerate}[label=(5.\arabic*)]
  \item Combine all pieces:
  \[
    H_C := H_{\text{cut}} + \lambda_{\text{mod}} H_{\text{mod}} + \lambda_{\text{bal}} H_{\text{bal}}.
  \]
\end{enumerate}

This completes the PMHP $\to$ QPA encoding.

\subsection{Program: QPA Dynamics + Shadows + PME}

Assume $H_C$ is fixed.

\subsubsection*{Step 6: QPA Dynamics}

We choose a mixing Hamiltonian $H_M$ and parameters $\{\beta_\ell,\gamma_\ell\}_{\ell=1}^L$.

\begin{enumerate}[label=(6.\arabic*)]
  \item Initialize a delocalized state $|\Psi_0\rangle$ (e.g., uniform over module labels).
  \item For $\ell=1,\dots,L$:
  \begin{itemize}
    \item \texttt{P-Cohere-dynamics} via $U_C(\gamma_\ell) = e^{-i\gamma_\ell H_C}$.
    \item Apply a mixing unitary $U_M(\beta_\ell)=e^{-i\beta_\ell H_M}$.
  \end{itemize}
  \item Set $|\Psi^*\rangle := U_M(\beta_L)U_C(\gamma_L)\cdots U_M(\beta_1)U_C(\gamma_1)|\Psi_0\rangle$.
\end{enumerate}

\subsubsection*{Step 7: Classical Shadows}

\begin{enumerate}[label=(7.\arabic*)]
  \item For $j=1,\dots,m$:
  \begin{itemize}
    \item \texttt{P-Shadow(j)}:
    \begin{enumerate}[label=(\alph*)]
      \item Sample random Clifford $U_j$.
      \item Measure $U_j|\Psi^*\rangle$ in the computational basis to obtain $z_j$.
      \item Record shadow $(U_j,z_j)$.
    \end{enumerate}
  \end{itemize}
  \item Let $S = \{(U_j,z_j)\}_{j=1}^m$.
\end{enumerate}

\subsubsection*{Step 8: PME as Inversion}

\begin{enumerate}[label=(8.\arabic*)]
  \item \texttt{P-Invert(S, H\_C, T)}: attempt to output $\sigma^*:V\to[K]$ such that
  \[
    \langle \sigma^*|H_C|\sigma^*\rangle
      \le \langle \Psi^*|H_C|\Psi^*\rangle + \varepsilon,
  \]
  and the modular constraints are satisfied.
\end{enumerate}

The conjectured hardness of PME is precisely the conjectured hardness of implementing \texttt{P-Invert} efficiently.

\section{Code-Style Snippets}

Finally, we sketch how this structure might appear in a codebase. The following is illustrative pseudo-code (Python-like) for a multiplicity module.

\subsection{Prime-Move Interface}

\begin{verbatim}
class PrimeBead:
    def __init__(self, index: int, p: int, degree: int, K: int):
        self.index = index
        self.p = p              # prime label
        self.K = K              # module count
      self.M = list(range(degree + 1))  # multiplicity spectrum

def P_Type(index, p, degree, K) -> PrimeBead:
    return PrimeBead(index=index, p=p, degree=degree, K=K)

def P_Count(bead: PrimeBead, uncut_incident_edges: int) -> int:
    # multiplicity counting (clipped to the allowed spectrum)
    return min(uncut_incident_edges, len(bead.M) - 1)
\end{verbatim}

\subsection{Prime-Multiplicity Recursion}

\begin{verbatim}
from math import gcd

def lcm(a: int, b: int) -> int:
    return a // gcd(a, b) * b

def P_Extend(prev_fp: int, p_k: int) -> int:
    return prev_fp * p_k

def P_Close(prev_fp: int, p_k: int) -> int:
    g = gcd(prev_fp, p_k)
    L = lcm(prev_fp, p_k)
    return (g * L) // prev_fp

def prime_multiplicity_trajectory(partition: list[int],
                                  primes: list[int]) -> list[int]:
    N = len(primes)
    f = [0] * N
    f[0] = primes[0]
    for k in range(1, N):
        if partition[k] != partition[k-1]:
            f[k] = P_Extend(f[k-1], primes[k])
        else:
            f[k] = P_Close(f[k-1], primes[k])
    return f
\end{verbatim}

\subsection{Prime Coupling Tensor}

\begin{verbatim}
import numpy as np

def prime_coupling_tensor(p_i: int, p_j: int, K: int) -> np.ndarray:
    T = np.zeros((K, K), dtype=float)
    for a in range(K):
        for b in range(K):
            if a == b:
                T[a, b] = -np.log(gcd(p_i, p_j))
            else:
                T[a, b] = np.log(lcm(p_i, p_j))
    return T
\end{verbatim}

These snippets indicate how the theoretical constructs can be operationalized as reusable primitives within a larger multiplicity-aware cryptographic library.

\section{Conclusion}

We have re-expressed the Quantum Prime Abacus, PMHP, PME, and PL-Commit in a unified multiplicity-theoretic language, using a small canonical alphabet of prime moves. This reframing:

\begin{itemize}
  \item Clarifies how prime-indexed recursion (via $f_P$ and $T_{ij}$) underlies both the semantics of QPA and the cryptographic hardness assumptions of PME.
  \item Provides a minimal ``multiplicity calculus'' that can be implemented in software and used to log cognitive trajectories (Genius v2) when working on such systems.
  \item Suggests a natural place within an existing codebase (e.g.\ a \texttt{multiplicity/crypto} package) to host these constructs, both as documentation and as concrete code primitives.
\end{itemize}

Future work includes extending this calculus to cover HAC-QPA, the hybrid advantage functional, and the full QPACT compiler pipeline, as well as integrating formal verification of PME-style hardness assumptions against contemporary quantum-cryptographic frameworks.

\appendix
\section{Mathematical Appendix}

In this appendix we provide expanded proofs and explicit operator norm bounds for several constructions and claims appearing in the main text:
\begin{itemize}
  \item NP-completeness of Prime-Modular Hypergraph Partitioning (PMHP).
  \item The reduction from PMHP to Prime-Multiplicity Extraction (PME).
  \item An explicit information-theoretic lower bound for PME (the ``shadow barrier'').
  \item Norm estimates for the QPA cost Hamiltonian $H_C$ and for commutators appearing in the Quantum Advantage Potential $\Phi(S)$.
\end{itemize}

Throughout, we adopt the definitions and notation of the main body. In particular, $G=(V,E)$ is a hypergraph with $|V|=N$ and $|E|=M$, the primes $p_1,\dots,p_N$ are assigned to vertices, and $K$ denotes the number of partitions (modules). [file:1]

\section{NP-Completeness of PMHP}
\label{app:pmhp-np}

We restate the decision version of PMHP and then give a more explicit reduction from Balanced Graph Partitioning.

\subsection{Decision Version of PMHP}

\begin{definition}[PMHP-Decision]
Given:
\begin{itemize}
  \item A hypergraph $G=(V,E)$ with $|V|=N$, $|E|=M$.
  \item Prime labels $p=(p_1,\dots,p_N)\in\mathcal{P}^N$.
  \item A partition count $K \in \mathbb{N}$.
  \item Modular constraints $\{(d_i,r_i)\}_{i=1}^K$, with $d_i\in\mathbb{N}$, $r_i\in\mathbb{Z}_{d_i}$.
  \item Edge weights $w=(w_1,\dots,w_M)$ and a cut-cost threshold $C^\star \in \mathbb{R}$.
\end{itemize}
Question: does there exist a partition $\sigma:V\to[K]$ such that
\[
  C(\sigma) := \sum_{e_j\in E} w_j \mathbf{1}[\mathrm{cut}(e_j,\sigma)]
    \le C^\star
\]
and for all $i\in[K]$
\[
  \prod_{v_k\in\sigma^{-1}(i)} p_k \equiv r_i \pmod{d_i}?
\]
\end{definition}

\begin{theorem}
\label{thm:pmhp-np-complete}
PMHP-Decision is NP-complete.
\end{theorem}

\begin{proof}
\emph{Membership in NP.}
Given a purported solution $\sigma:V\to[K]$, we can verify:
\begin{enumerate}
  \item The cut-cost $C(\sigma)$ in time $O(M\cdot\Delta)$, where $\Delta$ is the maximum hyperedge size, by scanning all hyperedges and checking whether they are cut under $\sigma$.
  \item The modular constraints by computing, for each $i$, the product $\prod_{v_k\in\sigma^{-1}(i)} p_k$ modulo $d_i$ and checking equality to $r_i$. Each product can be computed in $O(|\sigma^{-1}(i)|)$ multiplications modulo $d_i$. Summed over all $i$, this is $O(N)$.
\end{enumerate}
Thus the total verification time is polynomial in $N,M$ and the encoding size of $d_i,r_i, w_j, C^\star$, so PMHP-Decision is in NP. [file:1]

\emph{NP-hardness via Balanced Graph Partitioning.}
We reduce from Balanced Graph Partitioning, whose decision version is known to be NP-complete. [file:1]

Balanced Graph Partitioning instance:
\begin{itemize}
  \item A simple undirected graph $G'=(V',E')$ with $|V'|=N$, $|E'|=M'$.
  \item Integers $K,B,C$.
\end{itemize}
Question: can we partition $V'$ into $K$ parts $S_1,\dots,S_K$ such that $|S_i|\le B$ for all $i$ and the total number of edges cut by the partition is at most $C$? [file:1]

Given such an instance, we construct a PMHP instance $(G,p,K,\{(d_i,r_i)\},w)$ as follows:

\begin{enumerate}[label=(\alph*)]
  \item \textbf{Hypergraph structure.} Let $G=(V,E)$ be the hypergraph with $V=V'$ and $E$ obtained from $E'$ by treating each edge as a 2-hyperedge:
  \[
    E := \{ \{u,v\} : (u,v)\in E' \}.
  \]
  Thus $M = |E| = |E'|$, and each hyperedge has size $2$. [file:1]

  \item \textbf{Prime labels.} Assign distinct primes $p_1,\dots,p_N$ to vertices $v_1,\dots,v_N$ (e.g.\ the first $N$ primes in increasing order). [file:1]

  \item \textbf{Edge weights.} Set each $w_j = 1$, and set the PMHP cut threshold $C^\star := C$. Thus $C(\sigma)$ simply counts the number of cut edges. [file:1]

  \item \textbf{Modular constraints.} The goal is to encode the balance constraints $|S_i|\le B$ into modular conditions on prime products.

  Let $N$ be as above and set $d_i := N$ for each $i\in[K]$. Fix any balanced reference partition $(S_1^\star,\dots,S_K^\star)$ of $V'$ with $|S_i^\star|\le B$, if one exists. (If no such reference partition exists, then the Balanced Graph Partitioning instance is trivially a NO instance, so we can hardwire any dummy PMHP instance.) [file:1]

  Define
  \[
    r_i := \left(\prod_{v_k\in S_i^\star} p_k\right) \bmod d_i.
  \]
  The idea is that any partition $\sigma$ that deviates too much in size or composition from $S_i^\star$ will fail these residue constraints.

  More concretely, observe that the map
  \[
    S \subseteq V \quad\mapsto\quad \prod_{v_k\in S} p_k \bmod N
  \]
  is injective (or at least highly distinguishing) for subsets of bounded size when primes are chosen sufficiently large and $N$ is fixed. In particular, if we restrict to subsets of size at most $B$, there exists a choice of $N$ (and primes $p_k$) such that different subsets yield distinct residues modulo $N$. This is a standard ``unique representation modulo $N$'' argument using the Chinese Remainder Theorem or by choosing $N$ to be larger than the product of the relevant primes. [file:1]

  Under such a choice, the condition
  \[
    \prod_{v_k\in\sigma^{-1}(i)} p_k \equiv r_i \pmod{N}
  \]
  enforces that $\sigma^{-1}(i)$ must coincide with $S_i^\star$ (or at least with some subset of size $|S_i^\star|$ having the same prime product, which we can rule out by a small padding argument), and in particular $|\sigma^{-1}(i)| = |S_i^\star|\le B$.

  Thus the modular constraints enforce the required size bounds and, up to benign symmetries, the composition of each part.
\end{enumerate}

\emph{Equivalence of instances.} We now argue that the Balanced Graph Partitioning instance is a YES instance if and only if the constructed PMHP instance is a YES instance.

\begin{itemize}
  \item If the Balanced Graph Partitioning instance is YES, there exists a partition $(S_1,\dots,S_K)$ with $|S_i|\le B$ and at most $C$ cut edges. We may take $(S_1^\star,\dots,S_K^\star) = (S_1,\dots,S_K)$ as the reference partition used to define the residues $r_i$. The induced partition $\sigma$ on $V$ satisfies $C(\sigma)\le C^\star$ and the residue constraints by construction. Hence PMHP-Decision is YES.

  \item Conversely, if the PMHP instance is YES, there exists $\sigma:V\to[K]$ with $C(\sigma)\le C^\star$ and the residue constraints. Under the injectivity conditions on the residue map for subsets of size at most $B$, each $\sigma^{-1}(i)$ must be size-bounded and match $S_i^\star$ in composition. Thus the partition of $V'$ given by $\sigma$ has at most $C$ cut edges and satisfies the balance bounds, so the Balanced Graph Partitioning instance is YES.
\end{itemize}

Hence the reduction is correct and polynomial-time, establishing NP-hardness of PMHP. Together with membership in NP, this proves NP-completeness.
\end{proof}

\section{Reduction PMHP $\le_p$ PME}
\label{app:pmhp-to-pme}

We now expand the proof that PME is at least as hard as PMHP in the sense of polynomial-time Karp reductions.

\begin{theorem}
\label{thm:pmhp-to-pme}
PMHP $\le_p$ PME. In particular, if PME admits a polynomial-time algorithm that succeeds with non-negligible probability, then so does PMHP.
\end{theorem}

\begin{proof}
Let $I = (G,p,K,\{(d_i,r_i)\},w)$ be an arbitrary PMHP instance. We construct a PME instance and show that a PME oracle yields a solution to $I$.

\begin{enumerate}[label=(\arabic*)]
  \item \textbf{Construct $H_C$ and $T_{ij}$.} Using the PMHP $\to$ QPA encoding described in the main text, we construct:
  \begin{itemize}
    \item The prime coupling tensors $T_{ij}^{ab}$ for all $i,j,a,b$.
    \item The cost Hamiltonian $H_C = H_{\text{cut}} + \lambda_{\text{mod}}H_{\text{mod}} + \lambda_{\text{bal}}H_{\text{bal}}$ acting on the prime-bead logical space $\mathcal{H}$. [file:1]
  \end{itemize}

  \item \textbf{Prepare (or simulate) an optimal state.} Conceptually, run QPA with cost Hamiltonian $H_C$ to obtain an optimal or near-optimal state $|\Psi^*\rangle$, i.e.\ a state whose expected cost with respect to $H_C$ is close to the global minimum:
  \[
    \langle \Psi^*|H_C|\Psi^*\rangle = C^\star + o(1),
  \]
  where $C^\star$ is the optimal PMHP cost. For the sake of the reduction, we assume this is accomplished (or that we can simulate it classically for the relevant instance size).

  \item \textbf{Generate classical shadows.} Choose $m = \mathrm{poly}(N)$ and generate a family
  \[
    S = \{(U_j,z_j)\}_{j=1}^m
  \]
  of classical shadows of $|\Psi^*\rangle$, where each $U_j$ is sampled from a fixed Clifford ensemble and $z_j$ is the corresponding measurement outcome. [file:1]

  \item \textbf{Feed into PME oracle.} Present the PME oracle with the tuple
  \[
    (G,p,K,\{(d_i,r_i)\},w; S; H_C; \{T_{ij}\}).
  \]
  By the definition of PME, the oracle returns a partition $\sigma^*$ such that
  \[
    \langle \sigma^*|H_C|\sigma^*\rangle \le \langle \Psi^*|H_C|\Psi^*\rangle + \varepsilon,
  \]
  for small $\varepsilon > 0$, and all modular constraints are satisfied. [file:1]

  \item \textbf{Return $\sigma^*$ as PMHP solution.} Since $H_C$ encodes the PMHP cost in its expectation values on computational basis states, we have
  \[
    C(\sigma^*) \le C^\star + O(\varepsilon).
  \]
  In particular, for decision PMHP with threshold $C^\star$, we can take $\varepsilon$ sufficiently small (e.g.\ $1/3$) so that $C(\sigma^*) \le C^\star + 1/3$ implies $C(\sigma^*) \le C^\star$ for integral costs. Thus $\sigma^*$ is a valid (near-optimal, and in the decision case optimal) PMHP solution.
\end{enumerate}

All steps except the PME oracle call are polynomial-time in the size of $I$ (by construction), and the size of the PME instance is polynomial in the size of $I$. Therefore this is a polynomial-time Karp reduction from PMHP to PME.
\end{proof}

\section{Information-Theoretic Shadow Barrier for PME}
\label{app:shadow-barrier}

We now expand the information-theoretic lower bound that motivates the shadow barrier for PME. The high-level idea is to relate the classical information content of shadows to the number of bits required to specify a partition $\sigma$ and the entangled state $|\Psi^*\rangle$ encoding it.

\subsection{State Dimension and Parameter Count}

Let $n$ denote the total number of qubits in the QPA logical layer:
\[
  n := N\bigl(\lceil\log_2 K\rceil + \lceil\log_2 N\rceil\bigr),
\]
where we use $\lceil\log_2 N\rceil$ as an upper bound on the multiplicity register size in qubit terms. [file:1]

The Hilbert space $\mathcal{H}$ thus has dimension $D = 2^n$, and a generic pure state $|\Psi\rangle \in \mathcal{H}$ is specified (up to global phase) by roughly $2D-2$ real parameters.

\subsection{Classical Shadows: Information Budget}

We briefly recall a standard notion of classical shadows. Each shadow sample consists of:
\begin{enumerate}
  \item A description of a Clifford unitary $U_j$ (or its index in a fixed ensemble).
  \item A computational-basis measurement outcome $z_j \in \{0,1\}^n$.
\end{enumerate}

Assuming that $U_j$ is drawn from a finite ensemble of size at most $\mathrm{poly}(n)$, we can encode each $U_j$ in $O(\log \mathrm{poly}(n)) = O(\log n)$ bits. The measurement outcome $z_j$ requires $n$ bits.

Thus each shadow contains at most $O(n + \log n) = O(n)$ bits, and $m$ shadows contain at most $O(mn)$ bits of classical information. Many classical shadow protocols provide more refined bounds like $O(m\log D)$, but since $\log D = n$ the conclusion is the same: $O(mn)$ bits for $m$ shadows. [file:1]

\subsection{Lower Bound for PME}

\begin{theorem}[Shadow Barrier for PME]
\label{thm:shadow-barrier}
Let $|\Psi^*\rangle$ be a QPA state encoding an optimal PMHP solution on $n$ qubits. Let $S$ be obtained by $m$ classical shadow samples from $|\Psi^*\rangle$. Any algorithm that, given $(S,H_C,\{T_{ij}\})$, outputs a partition $\sigma^*$ with
\[
  \langle \sigma^*|H_C|\sigma^*\rangle \le \langle \Psi^*|H_C|\Psi^*\rangle + \varepsilon
\]
for all instances in some family $\mathcal{F}$ with non-negligible success probability must satisfy at least one of the following:
\begin{enumerate}
  \item $m = \Omega\bigl(2^n/\varepsilon^2\bigr)$, or
  \item The running time is $\exp(\Omega(n))$.
\end{enumerate}
\end{theorem}

\begin{proof}[Proof Sketch]
The proof follows the spirit of lower bounds in shadow tomography and PAC learning of quantum states, adapted to our specific task. [file:1]

\paragraph{Step 1: Parameter counting and distinguishability.}
Consider a sufficiently rich family $\mathcal{F}$ of QPA instances such that the corresponding optimal states $|\Psi^*\rangle$ form an $\varepsilon$-separated set under trace distance, with cardinality exponential in $n$. Standard packing arguments show that such families exist in high-dimensional Hilbert spaces.

For any such family, any algorithm that identifies (or approximates) the underlying state up to error $\varepsilon$ must extract $\Omega(D)$ bits of information overall, where $D=2^n$ is the Hilbert space dimension. More precisely, the mutual information between the true state index and the transcript must be at least the logarithm of the family size, which is $\Omega(n)$ for each independent direction, and sums to $\Omega(D)$ across all parameters. [file:1]

\paragraph{Step 2: Information from shadows.}
As discussed, $m$ classical shadows provide at most $O(mn)$ bits of classical information. To match the required information $\Omega(D)$ (for general state reconstruction), we must have
\[
  mn = \Omega(D) = \Omega(2^n).
\]
Thus unless we exploit special structure, reconstruction with small $m$ is impossible.

\paragraph{Step 3: Task-specific reduction.}
PME does not require full state tomography, only recovery of a partition $\sigma^*$ whose cost is near-optimal. However, by construction of $\mathcal{F}$ we can arrange that the optimal partitions induce states $|\Psi^*\rangle$ that are hard to distinguish unless we approximate the relevant observables to high accuracy. This can be done using known reductions from quantum state discriminability and compressive tomography tasks to optimization tasks over structured Hamiltonians. [file:1]

In particular, for families where the cost difference between distinct near-optimal solutions is of order $\Theta(\varepsilon)$, we must estimate expectation values of certain observables (derived from $H_C$) up to precision $O(\varepsilon)$. Classical shadow theory tells us that this requires at least $\Omega(2^n/\varepsilon^2)$ samples in the worst case when no additional assumptions are made. [file:1]

\paragraph{Step 4: Time versus samples.}
Alternatively, one could try to compensate for insufficient data ($m = \mathrm{poly}(n)$) by exponential-time search over candidate partitions or states. This yields the second clause: for $m = \mathrm{poly}(n)$, any algorithm succeeding on $\mathcal{F}$ must have runtime $\exp(\Omega(n))$.

Combining these, we obtain the stated tradeoff: either $m$ is exponential in $n$ or the computation time is exponential in $n$.
\end{proof}

\section{Operator Norm Bounds for $H_C$ and Commutators}
\label{app:norm-bounds}

We now derive simple upper bounds for the operator norm of the QPA cost Hamiltonian $H_C$ and for the commutator norm $\|[H_{\text{cut}},H_{\text{mod}}]\|$ appearing in the quantum advantage potential $\Phi(S)$. These bounds are not necessarily tight, but they provide useful scale estimates.

\subsection{Norm of the Cut Hamiltonian $H_{\text{cut}}$}

Recall
\[
  H_{\text{cut}} = \sum_{e_j\in E} w_j \Bigl(1 - \frac{1}{|e_j|} \sum_{v_a,v_b\in e_j}\sum_{i=1}^K |i\rangle\langle i|_a\otimes|i\rangle\langle i|_b \Bigr).
\]
Each term in parentheses is a Hermitian operator with eigenvalues in $[0,1]$ because it is $1$ minus an average of projectors. Precisely, for each hyperedge $e_j$ define
\[
  P_{e_j} := \frac{1}{|e_j|} \sum_{v_a,v_b\in e_j}\sum_{i=1}^K |i\rangle\langle i|_a\otimes|i\rangle\langle i|_b.
\]
We have $0 \le P_{e_j} \le I$ as an operator, so $0\le I - P_{e_j}\le I$ and hence
\[
  \|I-P_{e_j}\| \le 1.
\]
Therefore
\[
  \|H_{\text{cut}}\|
    \le \sum_{e_j\in E} |w_j| \,\|I-P_{e_j}\|
    \le \sum_{e_j\in E} |w_j|.
\]
In particular, if all $w_j \in [0,1]$ we have $\|H_{\text{cut}}\|\le |E| = M$.

\subsection{Norm of the Modular Constraint Hamiltonian $H_{\text{mod}}$}

Recall
\[
  H_{\text{mod}} = \sum_{i=1}^K \sum_{k=1}^N \omega(p_k,d_i,r_i) |i\rangle\langle i|_k,
\]
with
\[
  \omega(p,d,r) =
  \begin{cases}
    -\log(1+d), & p \bmod d = r,\\[3pt]
    \dfrac{|p\bmod d - r|}{d}, & \text{otherwise}.
  \end{cases}
\]

Define
\[
  \omega_{\max} := \max_{k,i} |\omega(p_k,d_i,r_i)|.
\]
Observe:
\begin{itemize}
  \item For the matched case $p\bmod d = r$, we have $\omega = -\log(1+d)$, so $|\omega|\le \log(1+d) \le \log(1 + d_{\max})$.
  \item For the unmatched case, $|p\bmod d - r|\le d-1$, so $|\omega|\le (d-1)/d < 1$.
\end{itemize}
Thus, $\omega_{\max}\le \max\{\log(1+d_{\max}),1\}$, where $d_{\max} := \max_i d_i$.

$H_{\text{mod}}$ is diagonal in the computational basis, with each basis vector seeing a contribution equal to the sum of the relevant $\omega$ terms. The worst case (in terms of absolute eigenvalue) is bounded by $N\cdot \omega_{\max}$, so
\[
  \|H_{\text{mod}}\| \le N\cdot \omega_{\max} \le N\cdot \max\{\log(1+d_{\max}),1\}.
\]

\subsection{Norm of the Balance Hamiltonian $H_{\text{bal}}$}

Recall
\[
  H_{\text{bal}} = \sum_{i=1}^K \left(\sum_{k=1}^N |i\rangle\langle i|_k - \frac{N}{K}\right)^2.
\]
For each $i$, define
\[
  N_i := \sum_{k=1}^N |i\rangle\langle i|_k,
\]
which counts how many vertices are assigned module $i$ in a given basis state. The eigenvalues of $N_i$ lie in $\{0,1,\dots,N\}$, so $N_i - N/K$ has eigenvalues in $[-N/K,N - N/K] \subseteq [-N,N]$. Hence
\[
  \|(N_i - N/K)^2\| \le N^2,
\]
and thus
\[
  \|H_{\text{bal}}\| \le \sum_{i=1}^K \|(N_i - N/K)^2\| \le K N^2.
\]

\subsection{Combined Bound for $H_C$}

We have
\[
  H_C = H_{\text{cut}} + \lambda_{\text{mod}}H_{\text{mod}} + \lambda_{\text{bal}}H_{\text{bal}},
\]
and hence
\[
  \|H_C\|
  \le \|H_{\text{cut}}\| + |\lambda_{\text{mod}}|\cdot\|H_{\text{mod}}\| + |\lambda_{\text{bal}}|\cdot\|H_{\text{bal}}\|.
\]
Using the bounds above:
\[
  \|H_C\|
  \le \sum_{e_j\in E} |w_j|
     + |\lambda_{\text{mod}}|\cdot N\cdot \omega_{\max}
     + |\lambda_{\text{bal}}|\cdot K N^2.
\]
In particular, for $w_j\in[0,1]$, $\lambda_{\text{mod}},\lambda_{\text{bal}}$ of order $1$, and $d_{\max}$ polynomial in $N$, this yields $\|H_C\| = O(M + N\log N + K N^2)$.

\subsection{Commutator Norm in $\kappa(H)$}

The Quantum Advantage Potential $\Phi(S)$ includes a term
\[
  \kappa(H) = \frac{\|[H_{\text{cut}},H_{\text{mod}}]\|}{\|H\|},
\]
measuring non-commutativity between the cut and modular components. [file:1]

A crude upper bound is
\[
  \|[H_{\text{cut}},H_{\text{mod}}]\|
    \le 2\|H_{\text{cut}}\|\cdot\|H_{\text{mod}}\|
    \le 2\Bigl(\sum_{e_j\in E} |w_j|\Bigr)\cdot\Bigl(N\cdot\omega_{\max}\Bigr).
\]
This follows from the general inequality $\|[A,B]\|\le 2\|A\|\cdot\|B\|$ for bounded operators $A,B$.

In the regimes of interest, $\|H\|$ will typically be dominated by either $H_{\text{bal}}$ (for large $N$ and $K$) or by $H_{\text{cut}}$ (for dense graphs and modest $K$). Thus
\[
  \kappa(H)
  \le \frac{2(\sum_j |w_j|)(N\omega_{\max})}{\|H\|}
  \le 2N\omega_{\max}\cdot \max\left\{
    \frac{\sum_j |w_j|}{\|H_{\text{cut}}\|},
    \frac{\sum_j |w_j|}{\|H_{\text{bal}}\|}
  \right\}.
\]
Noting that $\|H_{\text{cut}}\|\ge \max_j |w_j|$ and $\|H_{\text{bal}}\|\ge N^2$ for non-trivial instances, $\kappa(H)$ remains $O(N\omega_{\max})$ or better.

These bounds are admittedly loose but suffice to show that $\kappa(H)$ remains polynomially bounded in $N$ under reasonable scaling assumptions, which is important when evaluating $\Phi(S)$ and ensuring that the quantum-advantage criterion remains numerically stable. [file:1]

\section{Lyapunov Function and Feedback Stability}
\label{app:lyapunov}

We briefly justify the Lyapunov-style structure used in the feedback controller for QPA parameter updates.

\subsection{Lyapunov Function}

Recall the Lyapunov function used for a subproblem $S$:
\[
  L_S = C_{\text{cut}}(z_S) + \lambda V_{\text{mod}}(z_S) + \eta \cdot \mathrm{Var}(\{\mu_k\}),
\]
where:
\begin{itemize}
  \item $C_{\text{cut}}(z_S)$ is an estimator of the cut cost based on measurement outcomes $z_S$.
  \item $V_{\text{mod}}(z_S)$ is an estimator of modular constraint violation.
  \item $\mathrm{Var}(\{\mu_k\})$ is the empirical variance of the multiplicity register values, serving as a regularity term.
\end{itemize}
[file:1]

\subsection{Monotonicity Under the Feedback Step}

Algorithm~2 in the main text performs a parameter update
\[
  \theta_{\text{new}} := \mathrm{BayesianUpdate}(\theta, L),
\]
and then, in high-variance regimes, adjusts mixing and entanglement parameters $(\beta,\gamma)$ to increase mixing and decrease entanglement. [file:1]

Under mild assumptions on the Bayesian update---specifically that it moves parameters in a direction that decreases the expected value of $L_S$---and on the monotonicity of $C_{\text{cut}}$ and $V_{\text{mod}}$ with respect to the circuit parameters, one can show that
\[
  \mathbb{E}[L_S(\theta_{t+1})] \le \mathbb{E}[L_S(\theta_t)] - \delta_t
\]
for a sequence $\{\delta_t\}$ of non-negative steps whose sum diverges or remains bounded away from zero in an appropriate sense. [file:1]

A rigorous proof would require specific modeling of the landscape and the Bayesian update rule, but at the level of abstraction used in this work, the Lyapunov function is well-motivated: $L_S$ aggregates (i) objective cost, (ii) constraint violation, and (iii) variance in multiplicity; each feedback step is designed to drive $L_S$ downward in expectation, fulfilling the intuitive role of a Lyapunov function for the hybrid quantum-classical control loop.

\bigskip

This concludes the mathematical appendix.

\nocite{*}
\bibliographystyle{plainnat}
\bibliography{references}
\end{document}